\begin{textsample}{POS Dim 1 – df – Score 49.00 – df/df\_ef\_10.txt}  \label{ex:f1_pos_092}
LÍNGUA ESTRANGEIRA A Secretaria de Estado de Educação do Distrito Federal optou por assumir a perspectiva de Língua Estrangeira – LE neste Currículo para \textbf{que} esta matriz possibilite o ensino de qualquer idioma, destacando, porém, a obrigatoriedade do ensino da Língua Inglesa conforme estabelecido na Lei 13.415/2017.
Ressalta -se \textbf{que} todos os objetos de conhecimento e habilidades referentes a esse componente curricular na BNCC estão devidamente contemplados e ampliados como objetivos de aprendizagem e conteúdos, de forma a nortear o processo de \textbf{ensino-aprendizagem} da Língua Inglesa e de qualquer outra língua estrangeira.
Essa opção do Distrito Federal por \textbf{um} referencial amplo se \textbf{baseia} em sua vocação ao plurilinguismo no \textbf{que} \textbf{diz} respeito às políticas públicas voltadas ao ensino de línguas, especialmente no \textbf{que} se refere à oferta de variados idiomas nos Centros Interescolares de Línguas – CIL e aos projetos \textbf{que} visam possibilitar as aprendizagens de línguas em escolas regulares das redes de ensino.
Esse caráter plurilíngue encontra raízes profundas no processo histórico do ensino de línguas no país, \textbf{que} teve início com a chegada dos \textbf{Jesuítas}, em 1500.
Eles ensinavam a língua portuguesa, ainda como estrangeira, a povos \textbf{que} \textbf{aqui} se encontravam. \textbf{Uma} vez oficializada a língua portuguesa como língua nacional ( 1759 ), outras línguas começaram a ser introduzidas no sistema de ensino e a ampliar os currículos de língua estrangeira.
Nesse sentido, com a chegada da família real ( 1808 ), as línguas francesa e inglesa foram introduzidas no currículo escolar, e, alguns anos mais tarde, após a proclamação da República ( 1889 ), a língua inglesa e a alemã se tornaram opcionais no currículo.
No \textbf{século} XX, sucessivas reformas alternaram o ensino da língua estrangeira entre caráter obrigatório e optativo.
A Reforma Capanema ( 1942 ) fez do latim, do francês e do inglês componentes curriculares obrigatórios do então Colegial, \textbf{hoje} Ensino Médio, \textbf{enquanto} no Ginasial, \textbf{atual} Ensino Fundamental – Anos Finais, o latim foi substituído pelo espanhol ( cf.
REVISTA HELB, s.d., seção Linha do Tempo ).
Anos depois, a Lei de Diretrizes e Bases da Educação Nacional, Lei nº. 4.024 de 1961, retirou o caráter obrigatório do ensino de línguas ( BRASIL, 1961 ).
Em 1976, o Ministério da Educação – MEC resgatou parcialmente o ensino de línguas, tendo decretado a obrigatoriedade da Língua Estrangeira Moderna – LEM no então 2º grau ( Ensino Médio ).
Finalmente, a Lei de Diretrizes e Bases da Educação Nacional – LDB, Lei nº. 9.394 de 1996 ( BRASIL, 1996 ), tornou o ensino de LEM obrigatório desde os Anos Finais do Ensino Fundamental até o término da Educação Básica.
Entretanto, apesar de a mencionada Lei referir -se ao componente como Língua Estrangeira Moderna, o foco de sua oferta sempre tendeu a se \textbf{restringir} ao ensino da língua inglesa, com exceção de localidades brasileiras próximas de \textbf{fronteiras} ou marcadas por \textbf{forte} \textbf{influência} cultural estrangeira, fruto de processos migratórios.
Isso possibilitou a oferta de outras línguas de maior interesse dessas comunidades.
Além disso, políticas públicas de ensino de Língua Estrangeira – LE, como os CIL, inaugurados ainda na década de 1970 em Brasília, assim como experiências similares desenvolvidas em outras unidades da Federação com o fim de garantir acesso a cursos de línguas em instituições educacionais públicas, propiciaram o ensino especializado de outras línguas de forma gratuita.
Com o advento do Mercado Comum do Sul – Mercosul, organização intergovernamental fundada em 1991, o interesse pelo ensino do Espanhol cresceu e resultou na promulgação da Lei nº 11.161 de 2005, a qual tornou obrigatória a oferta da língua espanhola no Ensino Médio e facultativa nos currículos de Anos Finais ( BRASIL, 2005 ).
Na referida lei, determinou -se \textbf{que}, até 2010, todas as escolas de Ensino Médio brasileiras, públicas e privadas, deveriam ofertar o Espanhol, com acesso de caráter facultativo para os estudantes.
No entanto, a oferta dessa língua enfrentou \textbf{uma} série de desafios, como o fato da carga horária oferecida para o espanhol ser inferior à de língua inglesa e a existência de falhas no processo de formação inicial de professores.
De acordo com a LDB, em seu art. 26, os currículos do Ensino Fundamental e Médio deveriam contar com \textbf{uma} base nacional comum e \textbf{uma} parte diversificada.
Entretanto, a LDB sofreu mudanças estruturais com a publicação da Lei nº 13.415 de 16 de fevereiro de 2017, \textbf{que} estabeleceu o Inglês como oferta obrigatória, revogou a Lei nº 11.161/2005 e deixou a língua espanhola como opcional nos currículos do Ensino Médio.
Essas mudanças foram refletidas fortemente na BNCC, promulgada no final de 2017.
Assim, \textbf{um} dos primeiros objetivos deste referencial curricular está relacionado com a sensibilização para a linguagem, de modo a preparar o estudante para se posicionar positivamente em relação à diversidade \textbf{que} o estudo de LE apresenta, em \textbf{um} processo \textbf{que} empreende também compreender outras formas de estar e ser no mundo ( HEIDEGGER, 2005 ; KOTHE, 2013 ).
Tal sensibilização \textbf{diz} respeito à construção de atitudes valorativas frente a outras culturas, pontos de vista, maneiras de expressão e seres \textbf{humanos}.
Nesse contexto, a sensibilização para a linguagem requer, portanto, aprender com e sobre os indivíduos de \textbf{uma} sociedade e, nesse processo, aprender sobre si mesmo e a sociedade em \textbf{que} \textbf{vivemos}.
A linguagem e o modo como ela é usada marcam lugares sociais das pessoas, sejam elas jovens, idosas, mulheres ou \textbf{homens} ( cisgênero ou transgênero ), Lésbicas, Gays, Bissexuais, Travestis, Transexuais, Transgênero e Intersexuais – LGBTI, negros, brancos, indígenas, das periferias, do campo ou dos centros urbanos, entre outros.
É a abertura para aceitação de diferentes línguas e maneiras de estar e ser no mundo, em seu diálogo com os eixos transversais do currículo – Educação para a Diversidade ; Cidadania e Educação em e para os Direitos Humanos e Educação para a Sustentabilidade – \textbf{que} possibilita \textbf{uma} educação capaz de promover e fortalecer a formação de indivíduos autônomos, críticos, conscientes de si e acolhedores das diferenças e das dimensões humana e social de outras culturas e da sua própria.
Para isso, a organização da matriz de LE em objetivos e conteúdos se apoia em quatro eixos desenvolvimentais nos Anos Finais : Oralidade ( interação discursiva, compreensão oral, produção oral ) ; Compreensão e Produção Escrita ( estratégia de leitura, leitura e fruição, avaliação dos textos lidos, estratégias de escrita : escrita e pós-escrita, compreensão e produção escrita - prática de escritas ) ; Interculturalidade ( manifestação e evolução linguístico-culturais, comunicação intercultural ) e Práticas Mediadas pelas Tecnologias Digitais ( práticas investigativas, produção autoral e \textbf{partilha} de informações ).
Cada eixo está dividido em dimensões \textbf{que} ajudam a organizar o desenvolvimento do estudante e oferecem parâmetros para \textbf{que} professores possam acompanhar o processo de aprendizagem e/ou aquisição da língua por parte dos estudantes ( cf.
KRASHEN, 1982 ).
A oralidade, ainda \textbf{um} desafio em cursos de LE na Educação Básica ( FARIA, 2009 ), é marcada como \textbf{um} eixo orientador.
Esse eixo apresenta \textbf{uma} habilidade \textbf{que} tem sido deixada em segundo plano, quando não abandonada completamente, em relação à escrita.
Por esse eixo, espera -se \textbf{que} estudantes, com apoio de professores, possam sustentar interação verbal.
Para isso, é necessário \textbf{que} se desenvolvam as outras dimensões da oralidade : a compreensão e a produção orais.
O eixo Compreensão e Produção Escrita também apresenta \textbf{uma} via de mão dupla.
Ela \textbf{diz} respeito à capacidade de compreensão e interpretação de textos escritos e da expressão do pensamento por esse meio.
Por se tratar da modalidade escrita, é importante não se \textbf{perder} de vista \textbf{que} “ a aula de LE deve criar condições para \textbf{que} o \textbf{educando} possa engajar -se em atividades \textbf{que} \textbf{demandam} o uso da língua a partir de temáticas relevantes ” ( SCHLATER, 2009, s.p. ) para o estudante.
A inclusão do eixo Interculturalidade aponta para \textbf{uma} \textbf{abordagem} \textbf{contemporânea} \textbf{que} não condiz com \textbf{métodos} “ centrados nos aspectos formais da língua e nas amostras de linguagem descontextualizadas ” ( \textbf{SANTOS}, 2004, p. 94 ).
Para Santos ( 2004 ), essa é \textbf{uma} possibilidade de desenvolvimento temático no ensino de línguas, de modo a torná -lo mais significativo, desafiador e prazeroso para os agentes envolvidos no processo.
O eixo Práticas Mediadas Pelas Tecnologias Digitais \textbf{diz} respeito a \textbf{um} suporte \textbf{poderoso} de acesso e contato com a língua estudada.
As tecnologias digitais podem oferecer espaço para comunicação nas modalidades escrita e oral.
Elas são capazes de propiciar condições adequadas de produção e de recepção tanto na dimensão oral quanto na escrita da língua.
Outro fator relevante \textbf{que} as tecnologias permitem trazer como prática na sala de aula de língua estrangeira é a possibilidade de criação em diversos aspectos : produção de material audiovisual por parte dos estudantes, expressão escrita, acessos a materiais diversos ( cf. \textbf{SANTOS} ; BEATO ; ARAGÃO, s.d. ).
Trata -se, enfim, de \textbf{um} aliado com potencial para apoiar \textbf{uma} transformação nas práticas de ensino e aprendizagem na escola pública.
A consideração desses eixos em conjunto e em constante movimento, sem atitudes e procedimentos estanques e sem \textbf{fragmentações} de práticas de linguagem, possibilita a construção de \textbf{um} ensino flexível, pois aponta para \textbf{um} tipo modelar ( paradigmático ) de ensino, favorecendo a construção de propostas específicas em cada escola.
Nessa perspectiva, não é necessário foco especial em quaisquer dos elementos nos processos \textbf{que} levarão à aquisição de \textbf{uma} nova língua, pois quando o ensino e a aprendizagem de \textbf{uma} língua estrangeira são \textbf{vistos} como \textbf{um} processo complexo ( MORIN, 2003 ) e ecológico ( CAPRA, 1996 ), mais chances têm de cumprirem seu importante papel na educação para a cidadania e para os direitos humanos, para a sustentabilidade e para a diversidade, de modo a contemplar a natureza das diferenças de gênero, de intelectualidade, de raça/etnia, de orientação sexual, de pertencimento, de personalidade, de cultura, de patrimônio, de classe social, de diferenças motoras, sensoriais, enfim, a diversidade \textbf{vista} como possibilidade de adaptação e de sobrevivência como espécie na sociedade ( DISTRITO FEDERAL, 2014 ).
Da mesma forma, o ensino-aprendizagem de \textbf{uma} nova língua também pode contribuir para \textbf{uma} progressiva autonomia do \textbf{educando} \textbf{que} lhe propicie ampliação de sua perspectiva de mundo.
Isso é possível porque o principal eixo filosófico da matriz não se orienta por \textbf{uma} perspectiva estrutural/gramatical, o \textbf{que} permite a adoção de conteúdos relacionados a interesses e/ou necessidades de cada comunidade escolar.
Para circular na língua estudada ( MOURA, 2015 ), temáticas relativas à educação ambiental, consciência familiar, respeito ao próximo e valorização da própria identidade, assim como a apreciação de costumes e valores de outros povos, poderão contribuir para o desenvolvimento comunicacional do \textbf{aprendiz}.
Nessa perspectiva, é necessário cuidar para \textbf{que} não se valorizem apenas países ou culturas dominantes, de modo a desmitificar o ensino da língua, tornando -o mais crítico no \textbf{que} \textbf{diz} respeito a relações estabelecidas entre os povos ao longo de \textbf{séculos}.
A orientação crítica apontada acima e adotada neste Currículo se expressa, sobretudo, pela função social requerida do componente curricular no contexto em \textbf{que} se insere.
Nesse sentido, a ideia consiste em \textbf{que} o ensino da língua possa ser articulado com elementos da cultura, da história, da sociedade e das relações \textbf{que} se estabelecem no contexto do qual faz parte.
Na condução do trabalho pedagógico, o trajeto da língua deve ser considerado \textbf{atentando} -se para a \textbf{influência} da colonização e do patriarcado e das lutas descolonizadoras de resistência nas matrizes culturais \textbf{que} originam determinadas línguas.
Assim, cabe refletir sobre o lastro deixado pelo movimento de aculturação e apagamento a partir do domínio e extermínio de determinados povos ( como as populações negras e indígenas nos diversos continentes, entre outras ) ao se abordar a importância das línguas e culturas em contato/conflito.
Trazer tal debate à tona faz com \textbf{que} o ensino da língua não seja \textbf{uma} \textbf{mera} repetição de normas gramaticais ; ao contrário, transforma esse processo em aprendizagens críticas e reflexivas, \textbf{que} estão para além da reprodução mecânica de palavras ou termos.
Assim, atividades diversas e uso de recursos variados, tais como projetos, \textbf{tarefas}, conteúdos interdisciplinares, temas transversais, jogos, leitura, teatro, tecnologias, música, entre outros, em perspectiva comunicativa, servirão para propiciar ambiente temático adequado para o desenvolvimento das aprendizagens e \textbf{fomentar} a construção da autonomia dos estudantes no processo ( ALMEIDA FILHO, 1987 ). \textbf{Abordagens} \textbf{contemporâneas} de ensino de línguas apontam para a inclusão de dimensões \textbf{historicamente} \textbf{relegadas} à marginalidade dos processos de ensino centrados no caráter metalinguístico.
Essas dimensões incluem a \textbf{centralidade} no sentido, o desenvolvimento temático, o protagonismo estudantil e o desenvolvimento da pessoa.
Elas não \textbf{implicam}, entretanto, a \textbf{exclusão} de outras dimensões privilegiadas em \textbf{abordagens} \textbf{que} mantêm a língua como centro.
O redimensionamento dos elementos implícitos em propostas \textbf{contemporâneas} busca ajustes mais equilibrados tradicionalmente negados na \textbf{abordagem} majoritariamente instalada nos sistemas de ensino, a gramatical.
Assim, pode -se indicar \textbf{uma} \textbf{abordagem} de caráter comunicacional como parte do enfrentamento das dificuldades a caminho da superação dos resultados indesejados ainda registrados no ensino de LE no Brasil.
A matriz curricular de LE propõe \textbf{uma} alternativa aberta e plural na oferta linguística.
Esta proposta pressupõe o \textbf{compromisso} de diversos \textbf{atores} da educação para sua implementação : aponta para a necessidade de formação continuada institucionalmente estabelecida ; indica a necessidade de investimento nos diferentes níveis formais de apoio à implementação do currículo ; solicita reestruturação de elementos condicionantes da oferta de LE e conta com o \textbf{compromisso} profissional de professores na experimentação de novas \textbf{abordagens}.
Em relação ao papel do professor nesse processo, várias são as estratégias \textbf{que} poderão ser utilizadas por ele com fins de materialização de \textbf{uma} \textbf{abordagem} centrada no sentido.
Pode -se, por exemplo, instalar \textbf{uma} estratégia temática pela qual o planejamento de cursos e de aulas se dê a partir de itens temáticos, reservando a itens linguísticos o papel de apoio no programa.
Nessa \textbf{direção}, os eixos transversais podem \textbf{desencadear} o protagonismo estudantil, pois guardam potencial de promover reflexão crítica sobre processos de dominação entre povos, culturas ou classes sociais.
Outra possibilidade de trazer à \textbf{cena} \textbf{abordagens} \textbf{focadas} no sentido é o tratamento intercultural \textbf{que} se pode dar ao ensino de línguas, centrando -se ações e práticas na sensibilização cultural, no ensino de aspectos específicos e de tratamento de diferenças ou semelhanças culturais.
O potencial do uso de novas tecnologias nesse aspecto é enorme, já \textbf{que} pode estimular o \textbf{espírito} da pesquisa e propiciar o desenvolvimento da autonomia do estudante em seu processo de aprendizagem.
Outros esforços podem ser envidados no sentido de modificar o tradicional ensino centrado na estrutura da língua.
O Currículo \textbf{ora} apresentado disponibiliza oportunidade de reflexão sobre o ensino \textbf{contemporâneo} de LE e convida a experimentações de novas possibilidades, quiçá novos rumos e resultados nesse campo da educação.
A abertura do educador a essas possibilidades, apoiada por ações formativas \textbf{que} lhe propiciem o aprofundamento de conhecimentos, ampliação de suas concepções e explicitação e formalização da própria \textbf{abordagem} ( SANT’ANA, 2017 ), pode \textbf{configurar} -se como elemento-chave \textbf{que} permitirá a materialização dos objetivos de aprendizagem deste marco referencial.

% matched lemmas: abordagem, aprendiz, aqui, atentar, ator, atual, basear, cena, centralidade, compromisso, configurar, contemporâneo, demandar, desencadear, direção, dizer, educar, enquanto, ensino-aprendizagem, espírito, exclusão, focar, fomentar, forte, fragmentação, fronteira, historicamente, hoje, homem, humanos, implicar, influência, jesuíta, mero, método, ora, partilha, perder, poderoso, que, relegar, restringir, santo, século, tarefa, um, ver, vivar
\end{textsample}
